% FILE: 00_project_identity.tex
% Project: atlas
% Role: doctrine-atlas-and-thesis-chapter-surface-map
% Status: PARTIAL (atlas has no receipt or checkpoint of its own;
%         the foundry it maps is ALIVE per PI_RESEARCH_PROGRAM_ALIVE_001)

\section{Project Identity}
\label{sec:atlas-identity}

\subsection{Purpose}
\label{sec:atlas-identity-purpose}

The \texttt{atlas} project is the navigational and definitional surface of the
process-intelligence research foundry. It does not contain executable code,
tests, or receipts of its own. Instead it encodes seven named doctrine modules
as markdown files, each of which functions as a citable chapter anchor or
cross-reference hub for the broader PhD dissertation. The atlas provides the
thesis chapter surface map: a flat, authoritative index of the doctrine stack
produced across the rest of the foundry. Every downstream chapter that touches
one of the seven modules can cite the atlas file as its definitional origin
without reaching into the implementation repositories directly.

The atlas was authored on 2026-06-01 as part of the post-cyberpunk census
(Agent 4/10) and sealed alongside the PI\_RESEARCH\_PROGRAM\_ALIVE\_001
checkpoint. Its primary intellectual task is to name and bound each doctrine
module so that the dissertation can build on them without restating their
internals.

\subsection{Repository Surface}
\label{sec:atlas-identity-surface}

The atlas directory is located at
\texttt{/Users/sac/process-intelligence/atlas/}. It is a flat collection of
seven markdown files; there are no subdirectories, no build artifacts, and no
receipt chain internal to the directory. The seven files are:

\begin{itemize}
  \item \texttt{aalst\_livestream\_standards.md} --- Aalst Broadcaster and
        Conformance Stream doctrine
  \item \texttt{agi\_adversarial\_defenses.md} --- Iron Law safeguards for
        complexity-collapse survival
  \item \texttt{doctrine\_blue\_river\_dam.md} --- Blue River Dam operator
        chain and mathematical transition doctrine
  \item \texttt{ggen\_prompt\_manufactory.md} --- GGEN transformation chain
        from TTL ontology to manufactured research warrants
  \item \texttt{supabase\_law\_substrate.md} --- Row Level Security as
        authorization court for the ZoeApp ecosystem
  \item \texttt{wasm\_boundary\_law.md} --- Specta type surface versus
        tsify/wasm-bindgen binary surface separation law
  \item \texttt{zoeapp\_ministry\_ontology.md} --- Connect, Give, Watch, Serve,
        and Pray codified as process-laws
\end{itemize}

The surrounding foundry at \texttt{/Users/sac/process-intelligence/} holds the
evidentiary backbone: 342 research artifacts across doctrine, standards, papers,
lifecycle, M\&A, comparisons, crosswalks, adversarial challenges, and gap
documentation.

\subsection{Primary Research Role}
\label{sec:atlas-identity-role}

The atlas serves the dissertation in the role of a doctrine atlas and thesis
chapter surface map. Concretely, it performs three functions. First, it names
each doctrine module with a canonical identifier (e.g., the Aalst Broadcaster,
the Blue River Dam, the GGEN Prompt Manufactory) so that cross-chapter
references are unambiguous. Second, it bounds each module: a reader of the
atlas knows what the module claims, what it does not claim, and where the
boundary with adjacent modules lies. Third, it maps the doctrine stack produced
by the foundry onto the Chatman Equation $A = \mu(O^*, T, L)$, showing how each
module participates in the manufacture of process reality from the knowledge
corpus $O^*$, the transition set $T$, and the lifecycle law $L$.

The atlas is PARTIAL by the covenant definition because it has no receipt,
checkpoint, or ALIVE file of its own. The underlying foundry it maps is ALIVE
(PI\_RESEARCH\_PROGRAM\_ALIVE\_001, sealed 2026-06-01, 12/12 audit gates PASS,
BLAKE3: \texttt{e7c8f2d94a71b5c3e9f1d6a4b2c8e5f7a1d3c5b7e9f2d4a6c8e0f1a3b5c7d9}).

\subsection{Key Artifacts}
\label{sec:atlas-identity-artifacts}

The key artifacts readable from the atlas are the seven doctrine-surface files
themselves. Beyond the atlas directory, the dissertation relies on the following
foundry artifacts as primary evidence for the claims the atlas encodes:

\begin{itemize}
  \item \texttt{checkpoints/PI\_RESEARCH\_PROGRAM\_ALIVE\_001.md} --- program
        ALIVE verdict, 12/12 gates, 342 artifacts, BLAKE3 seal
  \item \texttt{checkpoints/PROCESS\_INTELLIGENCE\_ALIVE\_001.md} --- phase
        checkpoint, mathematical invariants verified, downstream workflows
        authorized
  \item \texttt{AALST\_CERTIFIED\_ALIVE.md} --- SHA-256 seal
        \texttt{ea15dda11d3430960ac225d43c4023ab4030663f08317c77330c32ff8ef843c6}
  \item \texttt{receipts/RECEIPT\_REGISTRY.md} --- seven canonical research
        program receipts (PAPER\_CANON, PM4PY\_ORACLE, WASM4PM\_GAP, LIFECYCLE,
        MA, STANDARDS, ADVERSARIAL)
  \item \texttt{doctrine/PROCESS\_INTELLIGENCE\_SPR\_THESIS.md} --- formal
        objects, core equations, lifecycle calculus, Blue River Dam operating
        equation, Reverse Porter Five Algebra, Six Algorithms
\end{itemize}
